\section{Source decompositions and Gauss operators}

\subsection{Normalization of the first source factor}

\begin{lemma}[The unsimplified first-cut contraction]
    \label{thm:mpug_raw_first_cut_contraction}
    \lean{MPOTensor.sourceCutM₁_adjointSimpleContraction_raw}
    \leanok
    \uses{def:mpu_source_matrix_one, def:mpu_source_weight,
        def:mpug_adjoint_simple_contraction}
    For an arbitrary tensor $A$ of physical dimension $d$ and bond
    dimension $D$, and an arbitrary $D\times D$ matrix $\rho$, put
    $M=M_1(A)$ and $W=I_d\otimes\rho$. Then
    \begin{align}
        M_1(s(A,\rho)) & =MM^\dagger W^{\mathsf T}M.
        \label{eq:mpug_raw_first_cut_contraction}
    \end{align}
    No positivity, canonical-form, or simplicity hypothesis is needed.
\end{lemma}
\begin{proof}\leanok
    \uses{def:mpu_source_matrix_one, def:mpu_source_weight,
        def:mpug_adjoint_simple_contraction}
    In the first-cut coordinates, multiplication and the physical
    Kronecker delta in $W^{\mathsf T}$ give
    \begin{align}
        (MM^\dagger W^{\mathsf T}M)_{(i,b),(a,j)}
         & =\sum_{p,q=0}^{d-1}\sum_{x,r,t=0}^{D-1}
        A^{i,p}_{xb}\,\overline{A^{q,p}_{xr}}\,
        A^{q,j}_{at}\,\rho_{tr}                    \\
         & =s(A,\rho)^{i,j}_{ab}.
    \end{align}
    The transpose preserves the boundary entry $\rho_{tr}$ in the
    defining contraction.
\end{proof}

\begin{lemma}[The first cut of the adjoint-simple contraction]
    \label{thm:mpug_first_cut_adjoint_contraction}
    \lean{MPOTensor.IsMPUCanonicalFormII.sourceCutM₁_adjointSimpleContraction}
    \leanok
    \discussion{7317}
    \uses{def:mpu_canonical_form_ii, def:mpu_weighted_source_factors,
        def:mpu_source_matrix_one, def:mpug_adjoint_simple_contraction}
    Let $\mathcal U$ be an MPU in canonical form II with recorded fixed
    matrix $\rho$, and let $X_1,Y_1$ be its chosen source factors for this
    same $\rho$. Construct $s$ from $\mathcal U,\rho$ as in
    Definition~\ref{def:mpug_adjoint_simple_contraction}. With the first
    cut convention $(M_1(A))_{(i,b),(a,j)}=A^{i,j}_{ab}$,
    \begin{align}
        M_1(s) & =X_1(Y_1Y_1^\dagger)Y_1.
        \label{eq:mpug_first_cut_adjoint_contraction}
    \end{align}
    In the source-factor expansion of
    \cite[lines~1326--1374]{FrancoRubio2025MPUGauging}, this factorization
    requires no simplicity assumption.
\end{lemma}

\begin{proof}\leanok
    \uses{thm:mpug_raw_first_cut_contraction,
        thm:mpu_source_factorizations, thm:mpu_source_factor_normalizations,
        def:mpu_weighted_source_factors, def:mpu_source_weight}
    With $W=I_d\otimes\rho$, diagonality gives $W^{\mathsf T}=W$.
    The cut factorization and weighted isometry then yield
    \begin{align}
        M_1(s)
         & = (X_1Y_1)(X_1Y_1)^\dagger W(X_1Y_1)       \\
         & =X_1(Y_1Y_1^\dagger)(X_1^\dagger W X_1)Y_1 \\
         & =X_1(Y_1Y_1^\dagger)Y_1.
    \end{align}
\end{proof}

\begin{lemma}[The first source Gram matrix is scalar]
    \label{thm:mpug_source_y_one_scalar_gram}
    \lean{MPOTensor.IsMPUCanonicalFormII.sourceY₁_mul_conjTranspose_eq_smul}
    \leanok
    \uses{def:mpu_canonical_form_ii, def:mpu_simple,
        def:mpdo_physical_adjoint, def:mpu_weighted_source_factors}
    Let $\mathcal U$ be an MPU in canonical form II, and suppose that
    both $\mathcal U$ and $\mathcal U^\dagger$ are simple. Construct the
    source factors from $\mathcal U$ and its recorded fixed matrix
    $\rho$. If $d_r$ is the first source-cut rank, then, as in the scalar
    Gram identity of \cite[lines~1375--1397]{FrancoRubio2025MPUGauging},
    there is a scalar $\delta\in\C$ such that
    \begin{align}
        Y_1Y_1^\dagger & =\delta I_{d_r}.
    \end{align}
\end{lemma}

\begin{proof}\leanok
    \uses{thm:mpug_adjoint_simple_scalar,
        thm:mpug_first_cut_adjoint_contraction,
        thm:mpu_source_factor_normalizations,
        thm:mpu_source_factor_right_inverses,
        thm:mpu_source_factorizations}
    By Lemma~\ref{thm:mpug_adjoint_simple_scalar}, $s=\delta\mathcal U$.
    Applying the first cut and
    Lemma~\ref{thm:mpug_first_cut_adjoint_contraction} gives
    \begin{align}
        X_1(Y_1Y_1^\dagger)Y_1 & =\delta X_1Y_1.
    \end{align}
    Multiplying on the left by $X_1^\dagger(I_d\otimes\rho)$ and on the
    right by $Z_1$ cancels the two outer factors via
    $X_1^\dagger(I_d\otimes\rho)X_1=I_{d_r}$ and $Y_1Z_1=I_{d_r}$.
\end{proof}

\begin{lemma}[Trace of the first source Gram matrix]
    \label{thm:mpug_source_y_one_gram_trace}
    \lean{MPOTensor.IsMPUCanonicalFormII.trace_sourceY₁_mul_conjTranspose}
    \leanok
    \uses{def:mpu_canonical_form_ii, def:mpu_simple,
        def:mpu_weighted_source_factors}
    Let $\mathcal U$ be a simple MPU in canonical form II, with physical
    dimension $d$, bond dimension $D$, first source-cut rank $d_r$, and
    recorded fixed matrix $\rho$. As in the physical-trace calculation of
    \cite[line~1398]{FrancoRubio2025MPUGauging}, which does not require
    simplicity of $\mathcal U^\dagger$, its chosen first source factor
    satisfies
    \begin{align}
        \operatorname{tr}(Y_1Y_1^\dagger) & =d.
    \end{align}
\end{lemma}

\begin{proof}\leanok
    \uses{thm:mpug_source_y_physical_contractions}
    Cyclic invariance gives
    $\operatorname{tr}(Y_1Y_1^\dagger)=\operatorname{tr}(Y_1^\dagger Y_1)$.
    Grouping the latter sum by its physical index gives
    \begin{align}
        \operatorname{tr}(Y_1^\dagger Y_1)
         & =\sum_{p=0}^{d-1}\sum_{r=0}^{d_r-1}\sum_{\alpha=0}^{D-1}
        \overline{(Y_1)_{r,(\alpha,p)}}(Y_1)_{r,(\alpha,p)}
        =\sum_{p=0}^{d-1}1=d
    \end{align}
    by Lemma~\ref{thm:mpug_source_y_physical_contractions}.
\end{proof}

\begin{lemma}[Normalization of the first source Gram matrix]
    \label{thm:mpug_source_y_one_normalization}
    \lean{MPOTensor.IsMPUCanonicalFormII.sourceY₁_mul_conjTranspose}
    \leanok
    \discussion{7317}
    \uses{def:mpu_canonical_form_ii, def:mpu_simple,
        def:mpdo_physical_adjoint, def:mpu_weighted_source_factors}
    Let $\mathcal U$ be an MPU in canonical form II of physical dimension
    $d$, and suppose that both $\mathcal U$ and $\mathcal U^\dagger$ are
    simple. Let $Y_1$ be the first source factor constructed from
    $\mathcal U$ and its recorded fixed matrix $\rho$, and let $d_r$ be
    the first source-cut rank. Then, as in the source-factor normalization
    proposition of \cite[lines~1064--1119]{FrancoRubio2025MPUGauging},
    \begin{align}
        Y_1Y_1^\dagger & =\frac{d}{d_r}I_{d_r}.
    \end{align}
    % Source: first identity of eq:MPUnice3. The other three constants
    % and the full prop:MPUsplus are not established by this entry.
\end{lemma}

\begin{proof}\leanok
    \uses{thm:mpug_source_y_one_scalar_gram,
        thm:mpug_source_y_one_gram_trace}
    Writing $Y_1Y_1^\dagger=\delta I_{d_r}$ by
    Lemma~\ref{thm:mpug_source_y_one_scalar_gram} and taking traces via
    Lemma~\ref{thm:mpug_source_y_one_gram_trace} gives $\delta d_r=d$.
    Since canonical form II gives $d>0$ and hence $d_r\ne0$, dividing as in
    \cite[line~1398]{FrancoRubio2025MPUGauging} yields $\delta=d/d_r$.
\end{proof}

\subsection{Normalization of the weighted second source factor}

\begin{lemma}[The first cut of the reflected contraction]
    \label{thm:mpug_adjoint_cut_second_gram}
    \lean{MPOTensor.IsMPUCanonicalFormII.sourceCutM₁_adjointSimpleContraction_physicalAdjoint}
    \leanok
    \discussion{7317}
    \uses{def:mpu_canonical_form_ii, def:mpu_weighted_source_factors,
        def:mpu_source_matrix_one, def:mpug_adjoint_simple_contraction,
        def:mpdo_physical_adjoint, def:mpu_source_weight}
    Let $\mathcal U$ be an MPU in canonical form II with recorded fixed
    matrix $\rho$, and let $X_2,Y_2$ be its chosen second source factors.
    Set $W=I_d\otimes\rho$ and let $s_\dagger$ be the contraction of
    Definition~\ref{def:mpug_adjoint_simple_contraction}, applied to
    $\mathcal U^\dagger$ and this same $\rho$. Then, in the reflected
    factorization for $Y_2(I_d\otimes\rho)Y_2^\dagger$ in the normalization
    proposition of \cite[lines~1064--1119 and 1398]{FrancoRubio2025MPUGauging},
    \begin{align}
        M_1(s_\dagger) & =Y_2^\dagger(Y_2WY_2^\dagger)X_2^\dagger,
    \end{align}
    without any simplicity hypothesis.
\end{lemma}

\begin{proof}\leanok
    \uses{thm:mpug_raw_first_cut_contraction,
        thm:mpu_physical_adjoint_source_cuts,
        thm:mpu_source_factorizations, thm:mpu_source_factor_normalizations,
        def:mpu_weighted_source_factors, def:mpu_source_weight}
    Lemma~\ref{thm:mpug_raw_first_cut_contraction} gives
    \begin{align}
        M_1(s(A,\rho)) & =M_1(A)M_1(A)^\dagger W^{\mathsf T}M_1(A).
    \end{align}
    The transpose records the boundary coefficient $\rho_{tr}$. Here
    $\rho$ is diagonal, so $W^{\mathsf T}=W$. Substituting
    $M_1(\mathcal U^\dagger)=M_2(\mathcal U)^\dagger=Y_2^\dagger X_2^\dagger$
    and canceling $X_2^\dagger X_2=I$ gives the factorization. All factors
    are those of $\mathcal U$, not newly chosen factors of its adjoint.
\end{proof}

\begin{lemma}[The weighted second source Gram matrix is scalar]
    \label{thm:mpug_source_y_two_scalar_gram}
    \lean{MPOTensor.IsMPUCanonicalFormII.sourceY₂_weighted_mul_conjTranspose_eq_smul}
    \leanok
    \uses{def:mpu_canonical_form_ii, def:mpu_simple,
        def:mpdo_physical_adjoint, def:mpu_weighted_source_factors,
        def:mpu_source_weight}
    Let $\mathcal U$ be an MPU in canonical form II with recorded fixed
    matrix $\rho$, and assume that both $\mathcal U$ and
    $\mathcal U^\dagger$ are simple. Let $Y_2$ be its chosen second
    source factor and $d_l$ its second source-cut rank. There is a scalar
    $\varepsilon\in\C$ such that
    \begin{align}
        Y_2(I_d\otimes\rho)Y_2^\dagger & =\varepsilon I_{d_l}.
    \end{align}
\end{lemma}

\begin{proof}\leanok
    \uses{def:mpug_adjoint_canonical_boundaries,
        thm:mpug_adjoint_simple_scalar, thm:mpu_physical_adjoint_identities,
        thm:mpug_adjoint_cut_second_gram,
        thm:mpu_physical_adjoint_source_cuts,
        thm:mpu_source_factorizations, thm:mpu_source_factor_normalizations,
        thm:mpu_source_factor_right_inverses}
    Passing canonical form II to $\mathcal U^\dagger$ with the same
    $\rho$, Lemma~\ref{thm:mpug_adjoint_simple_scalar} gives
    $s_\dagger=\varepsilon\mathcal U^\dagger$, since physical adjunction
    is involutive and both required simplicity hypotheses hold.
    Taking the first cut and applying
    Lemma~\ref{thm:mpug_adjoint_cut_second_gram} gives
    \begin{align}
        Y_2^\dagger(Y_2(I_d\otimes\rho)Y_2^\dagger)X_2^\dagger
         & =\varepsilon Y_2^\dagger X_2^\dagger.
    \end{align}
    Multiplying on the left by $Z_2^\dagger$ and on the right by $X_2$
    cancels the outer factors via $Z_2^\dagger Y_2^\dagger=I$ and
    $X_2^\dagger X_2=I$. The scalar comes from the adjoint contraction; it
    has not been identified with the scalar for the original tensor.
\end{proof}

\begin{lemma}[Trace of the weighted second source Gram matrix]
    \label{thm:mpug_source_y_two_gram_trace}
    \lean{MPOTensor.IsMPUCanonicalFormII.trace_sourceY₂_weighted_mul_conjTranspose}
    \leanok
    \uses{def:mpu_canonical_form_ii, def:mpu_simple,
        def:mpu_weighted_source_factors, def:mpu_source_weight}
    Let $\mathcal U$ be a simple MPU in canonical form II of physical
    dimension $d$, with recorded fixed matrix $\rho$ and chosen second
    source factor $Y_2$. Then, as in the physical-trace calculation for
    $Y_2$ in \cite[line~1398]{FrancoRubio2025MPUGauging}, which does not
    require simplicity of $\mathcal U^\dagger$,
    \begin{align}
        \operatorname{tr}\bigl(Y_2(I_d\otimes\rho)Y_2^\dagger\bigr) & =d.
    \end{align}
\end{lemma}

\begin{proof}\leanok
    \uses{thm:mpug_source_y_physical_contractions, def:mpu_source_weight}
    Cyclic invariance and expanding $I_d\otimes\rho$ give
    \begin{align}
        \operatorname{tr}\bigl(Y_2(I_d\otimes\rho)Y_2^\dagger\bigr)
         & =\sum_{p=0}^{d-1}\sum_{x,y=0}^{D-1}\sum_{l=0}^{d_l-1}
        \overline{(Y_2)_{l,(p,y)}}(Y_2)_{l,(p,x)}\rho_{xy}
        =\sum_{p=0}^{d-1}1=d
    \end{align}
    by Lemma~\ref{thm:mpug_source_y_physical_contractions} at $q=p$.
\end{proof}

\begin{lemma}[Normalization of the weighted second source Gram matrix]
    \label{thm:mpug_source_y_two_normalization}
    \lean{MPOTensor.IsMPUCanonicalFormII.sourceY₂_weighted_mul_conjTranspose}
    \leanok
    \discussion{7317}
    \uses{def:mpu_canonical_form_ii, def:mpu_simple,
        def:mpdo_physical_adjoint, def:mpu_weighted_source_factors,
        def:mpu_source_weight}
    Let $\mathcal U$ be an MPU in canonical form II of physical dimension
    $d$, and suppose that both $\mathcal U$ and $\mathcal U^\dagger$ are
    simple. Let $Y_2$ be the second source factor constructed from
    $\mathcal U$ and its recorded fixed matrix $\rho$, and let $d_l$ be
    the second source-cut rank. Then, as in the source-factor normalization
    proposition of \cite[lines~1064--1119]{FrancoRubio2025MPUGauging},
    \begin{align}
        Y_2(I_d\otimes\rho)Y_2^\dagger & =\frac{d}{d_l}I_{d_l}.
    \end{align}
    % Source: second identity of eq:MPUnice3. Neither physical X
    % normalization nor the full prop:MPUsplus follows from this entry alone.
\end{lemma}

\begin{proof}\leanok
    \uses{thm:mpug_source_y_two_scalar_gram,
        thm:mpug_source_y_two_gram_trace}
    Writing the Gram matrix as $\varepsilon I_{d_l}$ by
    Lemma~\ref{thm:mpug_source_y_two_scalar_gram}, its trace is $d$ by
    Lemma~\ref{thm:mpug_source_y_two_gram_trace}, so $\varepsilon d_l=d$.
    As $d>0$, the rank $d_l$ is nonzero, and the trace argument of
    \cite[line~1398]{FrancoRubio2025MPUGauging} gives $\varepsilon=d/d_l$.
\end{proof}

\subsection{Physical normalizations of the two source X factors}

\begin{lemma}[The canonical one-letter output contraction]
    \label{thm:mpug_one_letter_output_contraction}
    \lean{MPOTensor.IsMPUCanonicalFormII.oneLetter_output_contraction}
    \leanok
    \discussion{7317}
    \uses{def:mpu_canonical_form_ii, def:mpu_simple,
        def:mpdo_physical_adjoint}
    Let $\mathcal U$ be an MPU in canonical form II of physical dimension
    $d$ and bond dimension $D$, with recorded fixed matrix $\rho$.
    Suppose that $\mathcal U^\dagger$ is simple. For $0\leq i,j<d$,
    \begin{align}
        \sum_{x,\alpha,\beta=0}^{D-1}\sum_{p=0}^{d-1}
        \mathcal U^{i,p}_{x\alpha}
        \overline{\mathcal U^{j,p}_{x\beta}}\rho_{\beta\alpha}
         & =\delta_{i,j}.
    \end{align}
    In the reflected one-letter contraction used for the physical
    source-factor normalizations in
    \cite[lines~1120--1162 and 1398]{FrancoRubio2025MPUGauging}, the open
    indices are the physical output indices and simplicity of $\mathcal U$
    is not required.
\end{lemma}

\begin{proof}\leanok
    \uses{def:mpug_adjoint_canonical_boundaries,
        thm:mpug_one_letter_physical_contraction, def:mpdo_physical_adjoint}
    Applying Lemma~\ref{thm:mpug_one_letter_physical_contraction} to
    $\mathcal U^\dagger$ with its canonical presentation carrying the
    same $\rho$, where
    $(\mathcal U^\dagger)^{p,i}_{x\alpha}=\overline{\mathcal U^{i,p}_{x\alpha}}$,
    and taking double conjugates gives the displayed sum.
\end{proof}

\begin{lemma}[Physical normalization of the first source X factor]
    \label{thm:mpug_source_x_one_physical_normalization}
    \lean{MPOTensor.IsMPUCanonicalFormII.sourceX₁_physical_contraction}
    \leanok
    \discussion{7317}
    \uses{def:mpu_canonical_form_ii, def:mpu_simple,
        def:mpdo_physical_adjoint, def:mpu_weighted_source_factors}
    Let $\mathcal U$ be an MPU in canonical form II of physical dimension
    $d$ and bond dimension $D$. Assume that both $\mathcal U$ and
    $\mathcal U^\dagger$ are simple. Let $X_1$ be the first source factor
    constructed from $\mathcal U$ and its recorded fixed matrix $\rho$,
    and let $d_r$ be the first source-cut rank. For $0\leq i,j<d$,
    \begin{align}
        \sum_{r=0}^{d_r-1}\sum_{\alpha,\beta=0}^{D-1}
        (X_1)_{(i,\alpha),r}\rho_{\beta\alpha}
        \overline{(X_1)_{(j,\beta),r}}
         & =\frac{d_r}{d}\delta_{i,j}.
    \end{align}
    In this first physical source-factor normalization of
    \cite[lines~1120--1162]{FrancoRubio2025MPUGauging}, the matrix $\rho$
    acts on the virtual legs and the identity acts on the physical space.
\end{lemma}

\begin{proof}\leanok
    \uses{thm:mpug_source_y_one_normalization,
        thm:mpug_source_y_one_gram_trace, thm:mpug_one_letter_output_contraction,
        thm:mpu_source_factorizations, def:mpu_source_matrix_one}
    Factoring $M_1(\mathcal U)=X_1Y_1$ and applying
    Lemma~\ref{thm:mpug_source_y_one_normalization} yields
    \begin{align}
        M_1(\mathcal U)M_1(\mathcal U)^\dagger
         & =\frac{d}{d_r}X_1X_1^\dagger.
    \end{align}
    Multiplying the $((i,\alpha),(j,\beta))$ entry by
    $\rho_{\beta\alpha}$ and summing over $\alpha,\beta$ makes the left-hand
    side $\delta_{i,j}$ by Lemma~\ref{thm:mpug_one_letter_output_contraction}.
    The trace identity in Lemma~\ref{thm:mpug_source_y_one_gram_trace}
    gives $d_r\ne0$, and canonical form II gives $d>0$. Multiplication by
    $d_r/d$ gives the required physical contraction.
\end{proof}

\begin{lemma}[Physical normalization of the second source X factor]
    \label{thm:mpug_source_x_two_physical_normalization}
    \lean{MPOTensor.IsMPUCanonicalFormII.sourceX₂_physical_contraction}
    \leanok
    \discussion{7317}
    \uses{def:mpu_canonical_form_ii, def:mpu_simple,
        def:mpdo_physical_adjoint, def:mpu_weighted_source_factors}
    Let $\mathcal U$ be an MPU in canonical form II of physical dimension
    $d$ and bond dimension $D$. Assume that both $\mathcal U$ and
    $\mathcal U^\dagger$ are simple. Let $X_2$ be its chosen second source
    factor and let $d_l$ be the second source-cut rank. For $0\leq i,j<d$,
    the second physical source-factor normalization in
    \cite[lines~1120--1162]{FrancoRubio2025MPUGauging} holds:
    \begin{align}
        \sum_{l=0}^{d_l-1}\sum_{x=0}^{D-1}
        (X_2)_{(x,i),l}\overline{(X_2)_{(x,j),l}}
         & =\frac{d_l}{d}\delta_{i,j}.
    \end{align}
\end{lemma}

\begin{proof}\leanok
    \uses{thm:mpug_source_y_two_normalization,
        thm:mpug_source_y_two_gram_trace, thm:mpug_one_letter_output_contraction,
        thm:mpu_source_factorizations, def:mpu_source_matrix_two,
        def:mpu_source_weight}
    With $W=I_d\otimes\rho$, the factorization $M_2(\mathcal U)=X_2Y_2$
    and Lemma~\ref{thm:mpug_source_y_two_normalization} give
    \begin{align}
        M_2(\mathcal U)WM_2(\mathcal U)^\dagger
         & =\frac{d}{d_l}X_2X_2^\dagger.
    \end{align}
    Summing the $((x,i),(x,j))$ entry over $x$ and expanding the left-hand
    side with $\rho_{\alpha\beta}=\rho_{\beta\alpha}$ gives
    $\delta_{i,j}$ by Lemma~\ref{thm:mpug_one_letter_output_contraction}.
    The trace identity in Lemma~\ref{thm:mpug_source_y_two_gram_trace}
    gives $d_l\ne0$ with $d>0$, so multiplying by $d_l/d$ yields the result.
\end{proof}

Lemmas~\ref{thm:mpug_source_y_one_normalization},
\ref{thm:mpug_source_y_two_normalization},
\ref{thm:mpug_source_x_one_physical_normalization}, and
\ref{thm:mpug_source_x_two_physical_normalization} establish all four
identities of the source-factor normalization proposition in
\cite[lines~1064--1162]{FrancoRubio2025MPUGauging} for the compact source
factors under canonical form~II and simplicity of $\mathcal U$ and
$\mathcal U^\dagger$; no additional boundary or Gram identities are assumed.
% This establishes prop:MPUsplus for the constructed source factors.
% It does not identify arbitrary separately supplied factors or modify
% the scope of the earlier raw proportionality theorem.

\begin{lemma}[Comparison of factorizations from one-sided inverses]
    \label{thm:mpug_factorization_comparison}
    \lean{Matrix.factorization_comparison_of_one_sided_inverses}
    \leanok
    Let $A,B,\widetilde B,C$ be finite index sets, possibly empty, and let
    $\mathbb S$ be a semiring. Suppose that
    \begin{align}
        X            & \in\operatorname{Mat}_{A\times B}(\mathbb S),           \\
        Y            & \in\operatorname{Mat}_{B\times C}(\mathbb S),           \\
        \widetilde X & \in\operatorname{Mat}_{A\times\widetilde B}(\mathbb S), \\
        \widetilde Y & \in\operatorname{Mat}_{\widetilde B\times C}(\mathbb S)
    \end{align}
    satisfy $XY=\widetilde X\widetilde Y$. If
    $L\in\operatorname{Mat}_{\widetilde B\times A}(\mathbb S)$ and
    $R\in\operatorname{Mat}_{C\times B}(\mathbb S)$ satisfy
    $L\widetilde X=I_{\widetilde B}$ and $YR=I_B$, then the rectangular
    comparison matrix $K=LX\in\operatorname{Mat}_{\widetilde B\times B}(\mathbb S)$
    obeys
    \begin{align}
        K            & =\widetilde YR, \\
        X            & =\widetilde XK, \\
        \widetilde Y & =KY.
    \end{align}
    In these one-sided-inverse identities from the proof of
    \cite[lines~1052--1066]{FrancoRubio2025MPUGauging}, the intermediate
    index sets $B$ and $\widetilde B$ are independent, and no equality of
    their cardinalities is assumed or concluded.
\end{lemma}
\begin{proof}\leanok
    Associativity and the two inverse identities give
    \begin{align}
        LX & =(LX)(YR)=L(XY)R
        =L(\widetilde X\widetilde Y)R=\widetilde YR.
    \end{align}
    Multiplying $XY=\widetilde X\widetilde Y$ on the right by $R$ gives
    $X=\widetilde XK$. Multiplying it on the left by $L$ gives
    $\widetilde Y=KY$.
\end{proof}

\begin{lemma}[Reciprocal unitary gauges from a Kronecker comparison]
    \label{thm:mpug_reciprocal_unitary_gauges}
    \lean{Matrix.exists_reciprocal_unitary_gauges_of_comparison}
    \leanok
    For $i=1,2$, suppose that two compatible matrix factorizations are related
    by square comparison matrices $K_i$ through
    \begin{align}
        X_i & =\widetilde X_iK_i, & \widetilde Y_i & =K_iY_i.
        \label{eq:mpug_comparison_matrices}
    \end{align}
    If $K_1\otimes K_2$ is unitary, then there are a positive real number
    $\delta$ and unitary matrices $W_1,W_2$ such that
    \begin{align}
        \widetilde X_1 & =\delta^{-1}X_1W_1,          &
        \widetilde Y_1 & =\delta W_1^\dagger Y_1,       \\
        \widetilde X_2 & =\delta X_2W_2,              &
        \widetilde Y_2 & =\delta^{-1}W_2^\dagger Y_2.
        \label{eq:mpug_reciprocal_unitary_gauges}
    \end{align}
    Thus the positive scalar factors in the orientation of
    \cite[Lemma at lines~1052--1066]{FrancoRubio2025MPUGauging} are
    $\delta_1=\delta^{-1}$ and $\delta_2=\delta$, and hence
    $\delta_1\delta_2=1$.
    % Source anchor: arXiv:2502.20257, Lemma lem:deco, lines 1052--1066.
\end{lemma}

\begin{proof}\leanok
    \uses{thm:mpug_unitary_kronecker_factors}
    Applying Lemma~\ref{thm:mpug_unitary_kronecker_factors} to
    $K_1\otimes K_2$ gives $\delta>0$ and unitary matrices $U_1,U_2$ with
    $K_1=\delta U_1$ and $K_2=\delta^{-1}U_2$. Setting
    $W_i=U_i^\dagger$ and substituting into
    (\ref{eq:mpug_comparison_matrices}) gives
    (\ref{eq:mpug_reciprocal_unitary_gauges}), while positivity of $\delta$
    gives $\delta^{-1}\delta=1$.
\end{proof}

\subsection{Uniqueness of the source decompositions}

\begin{lemma}[Transformation of the source contraction]
    \label{thm:mpug_source_contraction_transport}
    \lean{MPOTensor.sourceU_contraction_transport}
    \leanok
    \discussion{7315}
    \uses{def:mpu_source_matrix_one, def:mpu_source_matrix_two}
    Let $U$ have physical dimension $d$ and bond dimension $D$, and write
    $r=\operatorname{rank}M_1(U)$ and $\ell=\operatorname{rank}M_2(U)$.
    Let $Y_1,\widetilde Y_1\in\operatorname{Mat}_{r\times(D\times d)}(\C)$
    and $Y_2,\widetilde Y_2\in\operatorname{Mat}_{\ell\times(d\times D)}(\C)$.
    For square matrices $K_1\in\operatorname{Mat}_r(\C)$ and
    $K_2\in\operatorname{Mat}_\ell(\C)$, suppose that
    $\widetilde Y_i=K_iY_i$ for $i=1,2$. Define
    $u,\widetilde u\colon\C^d\otimes\C^d\to\C^\ell\otimes\C^r$ by
    \begin{align}
        u_{(l,r'),(p,q)}
         & =\sum_{\beta=0}^{D-1}(Y_2)_{l,(p,\beta)}(Y_1)_{r',(\beta,q)}, \\
        \widetilde u_{(l,r'),(p,q)}
         & =\sum_{\beta=0}^{D-1}(\widetilde Y_2)_{l,(p,\beta)}
        (\widetilde Y_1)_{r',(\beta,q)}.
        \label{eq:mpug_raw_source_contractions}
    \end{align}
    Here $0\leq l<\ell$, $0\leq r'<r$, and $0\leq p,q<d$.
    Then $\widetilde u=(K_2\otimes K_1)u$.
\end{lemma}

\begin{proof}\leanok
    Substituting $\widetilde Y_i=K_iY_i$ into
    (\ref{eq:mpug_raw_source_contractions}) and interchanging the three finite
    sums isolates the coefficient of $u_{(l',r''),(p,q)}$ as
    $(K_2)_{l,l'}(K_1)_{r',r''}$, the corresponding entry of $K_2\otimes K_1$.
\end{proof}

\begin{lemma}[Reciprocal unitary gauges for the source decompositions]
    \label{thm:mpug_source_decomposition_gauges}
    \lean{MPOTensor.exists_reciprocal_unitary_source_gauges}
    \leanok
    \discussion{7315}
    \uses{def:mpu_source_matrix_one, def:mpu_source_matrix_two}
    Let $U$ have physical dimension $d$ and bond dimension $D$, with
    $r=\operatorname{rank}M_1(U)$ and $\ell=\operatorname{rank}M_2(U)$.
    Suppose that
    \begin{align}
        X_1,\widetilde X_1 & \in\operatorname{Mat}_{(d\times D)\times r}(\C),   &
        Y_1,\widetilde Y_1 & \in\operatorname{Mat}_{r\times(D\times d)}(\C),      \\
        X_2,\widetilde X_2 & \in\operatorname{Mat}_{(D\times d)\times\ell}(\C), &
        Y_2,\widetilde Y_2 & \in\operatorname{Mat}_{\ell\times(d\times D)}(\C)
    \end{align}
    satisfy $X_iY_i=\widetilde X_i\widetilde Y_i=M_i(U)$ for $i=1,2$.
    Suppose there are matrices $L_i,R_i$ with
    $L_i\widetilde X_i=I$ and $Y_iR_i=I$.
    Assume that the maps $u,\widetilde u\colon\C^d\otimes\C^d\to
        \C^\ell\otimes\C^r$ with coordinates
    \begin{align}
        u_{(l,r'),(p,q)}
         & =\sum_{\beta=0}^{D-1}(Y_2)_{l,(p,\beta)}(Y_1)_{r',(\beta,q)}, \\
        \widetilde u_{(l,r'),(p,q)}
         & =\sum_{\beta=0}^{D-1}(\widetilde Y_2)_{l,(p,\beta)}
        (\widetilde Y_1)_{r',(\beta,q)}
    \end{align}
    are unitary. The indices range over $0\leq l<\ell$, $0\leq r'<r$,
    and $0\leq p,q<d$. Then there exist $\delta>0$ and unitaries
    $W_1\in\operatorname{Mat}_r(\C)$, $W_2\in\operatorname{Mat}_\ell(\C)$ such that
    \begin{align}
        \widetilde X_1 & =\delta^{-1}X_1W_1,          &
        \widetilde Y_1 & =\delta W_1^\dagger Y_1,       \\
        \widetilde X_2 & =\delta X_2W_2,              &
        \widetilde Y_2 & =\delta^{-1}W_2^\dagger Y_2.
        \label{eq:mpug_source_decomposition_gauges}
    \end{align}
    Thus $\delta_1=\delta^{-1}$ and $\delta_2=\delta$ satisfy
    $\delta_1\delta_2=1$, establishing the decomposition lemma of
    \cite[lines~1052--1066]{FrancoRubio2025MPUGauging} from the cut factorizations,
    one-sided inverses, and unitarity of $u,\widetilde u$ alone, without requiring
    unitarity of the second contraction $v$.
    % Source lem:deco, issue #7315. No clause (c), weighted normalization,
    % canonical-form-II, density, or physical-adjoint hypothesis is used.
    % This result does not assert completion of the separate #7313 assembly.
\end{lemma}

\begin{proof}\leanok
    \uses{thm:mpug_factorization_comparison,
        thm:mpug_source_contraction_transport, thm:mpug_reciprocal_unitary_gauges}
    If $d=0$, every factor equality has an empty physical row or column
    index. Thus $\delta=1$ and identity gauges satisfy all four
    equalities, and $\delta^{-1}\delta=1$.
    Suppose now that $d>0$. Since $u$ is an isometry,
    $\ell r\geq d^2>0$; hence both
    source ranks are positive. Lemma~\ref{thm:mpug_factorization_comparison}
    gives $K_i=L_iX_i=\widetilde Y_iR_i$ with
    $X_i=\widetilde X_iK_i$ and $\widetilde Y_i=K_iY_i$.
    Lemma~\ref{thm:mpug_source_contraction_transport} yields
    $\widetilde u=(K_2\otimes K_1)u$. Therefore
    $K_2\otimes K_1=\widetilde u u^\dagger$ is unitary.
    Interchanging the two tensor factors makes $K_1\otimes K_2$ unitary, and
    Lemma~\ref{thm:mpug_reciprocal_unitary_gauges} then gives
    (\ref{eq:mpug_source_decomposition_gauges}). In particular, if
    $K_1=\delta V_1$ and $K_2=\delta^{-1}V_2$ with $V_i$ unitary,
    the required orientation is $W_i=V_i^\dagger$.
\end{proof}

\begin{corollary}[Unitary gauges for normalized source decompositions]
    \label{cor:mpug_normalized_source_decomposition_gauges}
    \lean{MPOTensor.exists_unitary_source_gauges_of_isometric_second_factors}
    \leanok
    \uses{thm:mpug_source_decomposition_gauges}
    Under the hypotheses of Lemma~\ref{thm:mpug_source_decomposition_gauges},
    suppose in addition that
    \begin{align}
        X_2^\dagger X_2 & =I_\ell,
                        & \widetilde X_2^\dagger\widetilde X_2 & =I_\ell.
    \end{align}
    In this isometrically normalized specialization of
    \cite[lines~1052--1066]{FrancoRubio2025MPUGauging}, the additional
    isometry hypotheses remove the scalar freedom: there exist unitaries
    $W_1\in\operatorname{Mat}_r(\C)$ and
    $W_2\in\operatorname{Mat}_\ell(\C)$ such that, for $i=1,2$,
    \begin{align}
        \widetilde X_i & =X_iW_i,
                       & \widetilde Y_i & =W_i^\dagger Y_i.
    \end{align}
\end{corollary}

\begin{proof}\leanok
    \uses{thm:mpug_source_decomposition_gauges}
    If $d=0$, all four equalities have an empty physical index and hold
    with identity gauges. If $d>0$, unitarity of $u$ gives $r\ell=d^2>0$,
    so $\ell>0$. Lemma~\ref{thm:mpug_source_decomposition_gauges} gives
    $\widetilde X_2=\delta X_2W_2$ with $\delta>0$. The two isometry
    identities and unitarity of $W_2$ imply
    $I_\ell=\delta^2I_\ell$. Since $\ell>0$, a diagonal entry gives
    $\delta^2=1$, hence $\delta=1$. Substitution into the four comparison
    identities proves the claim.
\end{proof}

\begin{definition}[Unitary factorization comparison]
    \label{def:mpug_unitary_factorization_comparison}
    \lean{Matrix.unitaryFactorizationComparison,
        Matrix.unitaryFactorizationComparison_coe}
    \leanok
    Let $I$ be a set and let $R_{a,b}$ be a family of unitary matrices of a
    fixed size, indexed by $(a,b)\in I^2$. Following the orientation in the
    corollary after Proposition~3
    of~\cite[lines~2782--2821]{FrancoRubio2025MPUGauging}, the unitary
    comparison from $(a,b)$ to the reference pair $(c,d)$ is
    \begin{align}
        \lambda_{a,b}^{c,d}
         & =(R_{c,d})^{-1}R_{a,b}
        =R_{c,d}^\dagger R_{a,b}.
        \label{eq:mpug_unitary_factorization_comparison}
    \end{align}
    % Source anchor: arXiv:2502.20257, Corollary after Proposition 3,
    % lines 2782--2821, especially the defining formula at line 2820.
    In that corollary, $I=G$ and comparisons are used when $ab=cd$; this
    equality selects the factorizations and is not needed in the unitary
    calculation.
\end{definition}

\begin{lemma}[Identity, inverse, composition, and right-reference laws]
    \label{thm:mpug_unitary_factorization_comparison_laws}
    \lean{Matrix.unitaryFactorizationComparison_self,
        Matrix.unitaryFactorizationComparison_inv,
        Matrix.unitaryFactorizationComparison_trans,
        Matrix.unitaryFactorizationComparison_of_reference_eq_one,
        Matrix.unitaryFactorizationComparison_right_reference}
    \leanok
    \uses{def:mpug_unitary_factorization_comparison}
    For all indices $a,b,c,d,f,g\in I$,
    \begin{align}
        \lambda_{a,b}^{a,b} & =\mathbf 1.
    \end{align}
    Reversing the comparison gives
    \begin{align}
        \left(\lambda_{a,b}^{c,d}\right)^{-1}
         & =\lambda_{c,d}^{a,b}.
    \end{align}
    Comparisons compose as
    \begin{align}
        \lambda_{a,b}^{f,g}
         & =\lambda_{c,d}^{f,g}\lambda_{a,b}^{c,d}.
        \label{eq:mpug_unitary_factorization_comparison_laws}
    \end{align}
    If $R_{c,d}=\mathbf 1$, then
    $\lambda_{a,b}^{c,d}=R_{a,b}$. In particular, if $I=G$ is a group, the
    normalization $R_{e,x}=\mathbf 1$ gives the right-reference specialization
    $\lambda^R_{g,h}=\lambda_{g,h}^{e,gh}$ in the cited corollary,
    \begin{align}
        R_{g,h} & =\lambda_{g,h}^{e,gh}.
        \label{eq:mpug_unitary_factorization_right_reference}
    \end{align}
    % Source anchor: arXiv:2502.20257, Corollary after Proposition 3,
    % line 2817 (right reference and composition).
\end{lemma}

\begin{proof}\leanok
    \uses{def:mpug_unitary_factorization_comparison}
    The identity, inverse, and right-reference formulas follow from
    $R_{a,b}^{-1}R_{a,b}=\mathbf 1$, inversion of a product, and
    $\mathbf 1^{-1}R_{a,b}=R_{a,b}$. For composition, inserting the identity at
    the intermediate pair gives
    \begin{align}
        (R_{f,g})^{-1}R_{a,b}
         & =(R_{f,g})^{-1}R_{c,d}(R_{c,d})^{-1}R_{a,b}.
    \end{align}
\end{proof}

\subsection{Gauss operators and Gauss projectors}

\begin{definition}[Gauss action on the gauge legs]
    \label{def:mpug_gauss_leg_action}
    \lean{TNLean.Algebra.gaussLegAction,
        TNLean.Algebra.gaussLegAction_apply}
    \leanok
    \discussion{7503}
    Let $G$ be a group. Following the two gauge-leg labels in the local Gauss
    operator of \cite[lines~3380--3387]{FrancoRubio2025MPUGauging}, for
    $g\in G$ define the permutation $T_g\colon G^2\to G^2$ by
    \begin{align}
        T_g(a,b) & =\left(ag^{-1},gb\right).
        \label{eq:mpug_gauss_leg_action}
    \end{align}
\end{definition}

\begin{lemma}[Composition of the Gauss leg action]
    \label{thm:mpug_gauss_leg_action_mul}
    \lean{TNLean.Algebra.gaussLegAction_one,
        TNLean.Algebra.gaussLegAction_mul}
    \leanok
    \discussion{7503}
    \uses{def:mpug_gauss_leg_action}
    For all $g,h\in G$,
    \begin{align}
        T_e          & =\id,    \\
        T_g\circ T_h & =T_{gh}.
        \label{eq:mpug_gauss_leg_action_mul}
    \end{align}
\end{lemma}

\begin{proof}\leanok
    \uses{def:mpug_gauss_leg_action}
    The two coordinates give
    \begin{align}
        T_g(T_h(a,b))
         & =\left(ah^{-1}g^{-1},ghb\right)
        =\left(a(gh)^{-1},(gh)b\right).
    \end{align}
\end{proof}

\begin{definition}[Local Gauss operator]
    \label{def:mpug_gauss_operator}
    \lean{TNLean.Algebra.gaussOperator,
        TNLean.Algebra.gaussOperator_apply_target,
        TNLean.Algebra.gaussOperator_apply_of_ne}
    \leanok
    \discussion{7503}
    \uses{def:mpug_unitary_factorization_comparison,
        def:mpug_gauss_leg_action}
    Let $G$ be finite, let $R_{a,b}\in\mathrm U(n)$ for $a,b\in G$, and take
    the basis $\ket{i,a,b}$ of $\C^n\otimes\C[G]^{\otimes2}$. As in
    \cite[lines~3380--3387]{FrancoRubio2025MPUGauging}, define the local Gauss
    operator
    \begin{align}
        \mathcal G_g
         & =\sum_{a,b\in G}
        \lambda_{a,b}^{ag^{-1},gb}
        \otimes\ket{ag^{-1},gb}\!\bra{a,b}.
        \label{eq:mpug_gauss_operator}
    \end{align}
    Equivalently, the block from $(a,b)$ to $T_g(a,b)$ is
    $(R_{T_g(a,b)})^{-1}R_{a,b}$ and every other block in that column
    vanishes.
\end{definition}

\begin{lemma}[Gauss operators form a unitary representation]
    \label{thm:mpug_gauss_representation}
    \lean{TNLean.Algebra.gaussOperator_one,
        TNLean.Algebra.gaussOperator_mul,
        TNLean.Algebra.star_gaussOperator,
        TNLean.Algebra.gaussOperator_mem_unitaryGroup,
        TNLean.Algebra.gaussRepresentation,
        TNLean.Algebra.gaussRepresentation_coe}
    \leanok
    \discussion{7503}
    \uses{def:mpug_gauss_operator}
    The local Gauss operators satisfy
    \begin{align}
        \mathcal G_e & =\id.
    \end{align}
    Their product law is
    \begin{align}
        \mathcal G_g\mathcal G_h & =\mathcal G_{gh}.
        \label{eq:mpug_gauss_representation}
    \end{align}
    Finally,
    \begin{align}
        \mathcal G_g^\dagger & =\mathcal G_{g^{-1}}.
    \end{align}
    Hence $g\mapsto\mathcal G_g$ is a unitary representation, establishing the
    representation part of the local Gauss-law proposition in
    \cite[lines~3380--3387]{FrancoRubio2025MPUGauging}; invariance of the
    gauged tensor is a separate assertion and is not included here.
\end{lemma}

\begin{proof}\leanok
    \uses{thm:mpug_unitary_factorization_comparison_laws,
        thm:mpug_gauss_leg_action_mul,
        def:mpug_gauss_operator}
    In a column labelled by $(a,b)$, multiplication has only the intermediate
    gauge label $T_h(a,b)$. The surviving matter block is
    \begin{align}
        \lambda_{T_h(a,b)}^{T_g(T_h(a,b))}
        \lambda_{a,b}^{T_h(a,b)}
         & =\lambda_{a,b}^{T_{gh}(a,b)},
    \end{align}
    by the comparison composition law and
    $T_gT_h=T_{gh}$. At $e$, every surviving comparison is the identity.
    Reversing a comparison gives
    $(\lambda_{a,b}^{c,d})^\dagger=\lambda_{c,d}^{a,b}$, so taking the adjoint
    gives $\mathcal G_g^\dagger=\mathcal G_{g^{-1}}$. The product with this
    adjoint is $\mathcal G_e=\id$, which proves unitarity.
\end{proof}

\begin{definition}[Unplaced local Gauss projector kernel]
    \label{def:mpug_gauss_projector}
    \lean{TNLean.Algebra.gaussProjector,
        TNLean.Algebra.gaussProjector_eq_average}
    \leanok
    \discussion{7578}
    \uses{def:mpug_finite_group_unitary_average,
        thm:mpug_gauss_representation}
    Let $G$ be finite and let $\mathcal G_g$ act on
    $\mathcal H=\C^n\otimes\C[G]^{\otimes2}$. Define the unplaced local
    projector kernel
    \begin{align}
        P_G(R) & =\frac{1}{|G|}\sum_{g\in G}\mathcal G_g.
        \label{eq:mpug_gauss_projector}
    \end{align}
    This is the local average introduced in
    \cite[lines~3494--3499]{FrancoRubio2025MPUGauging}, before a position on
    the chain is chosen. This definition neither places the kernel on
    $(j,j+1)$ nor asserts invariance of the gauged matrix product state.
\end{definition}

\begin{lemma}[The local Gauss average is a projection]
    \label{thm:mpug_gauss_projector_projection}
    \lean{TNLean.Algebra.isStarProjection_gaussProjector}
    \leanok
    \discussion{7578}
    \uses{def:mpug_gauss_projector}
    The kernel in~(\ref{eq:mpug_gauss_projector}) satisfies
    \begin{align}
        P_G(R)^2       & =P_G(R), \\
        P_G(R)^\dagger & =P_G(R).
    \end{align}
\end{lemma}

\begin{proof}\leanok
    \uses{thm:mpug_finite_group_unitary_average_projection,
        thm:mpug_gauss_representation}
    Identifying the matrix average with the averaging operator of the unitary
    representation $g\mapsto\mathcal G_g$ makes this operator idempotent by
    the invariant-space averaging theorem. For self-adjointness,
    $\mathcal G_g^\dagger=\mathcal G_{g^{-1}}$ and reindexing the finite sum by
    the bijection $g\mapsto g^{-1}$ yield $P_G(R)^\dagger=P_G(R)$.
\end{proof}

\begin{definition}[Gauss projector on a periodic chain]
    \label{def:mpug_placed_gauss_projector}
    \lean{MPOTensor.gaussLocalCoordinateEquiv,
        MPOTensor.gaussLocalCoordinateEquiv_apply_zero,
        MPOTensor.gaussLocalCoordinateEquiv_apply_one,
        MPOTensor.gaussWindowProjector,
        MPOTensor.placedGaussProjector}
    \leanok
    \discussion{7605}
    \uses{def:mpug_gauss_projector}
    Let the matter degree of freedom at one site have dimension $d$.  The
    local kernel has row and column coordinates
    \begin{align}
        \bigl((i_0,i_1),(a,b)\bigr)
         & \in (\{0,\ldots,d-1\}^{\{0,1\}})\times G^2.
    \end{align}
    Identifying these with two site coordinates by the bijection
    \begin{align}
        \kappa\bigl((i_0,i_1),(a,b)\bigr)
         & =\bigl((i_0,a),(i_1,b)\bigr),
    \end{align}
    followed at each site by a fixed enumeration of
    $\{0,\ldots,d-1\}\times G$, let $\widetilde P_G(R)$ denote the matrix
    obtained from $P_G(R)$ by this simultaneous row and column relabelling.
    Following the position-dependent projector $P_j$
    in~\cite[lines~3494--3499]{FrancoRubio2025MPUGauging}, define its placement
    on a periodic chain of length $N\geq2$ by
    \begin{align}
        P_j(R) & =\operatorname{Embed}_{j,j+1}\bigl(\widetilde P_G(R)\bigr),
        \label{eq:mpug_placed_gauss_projector}
    \end{align}
    where the second site is $j+1$ modulo $N$.
    % Source anchor: arXiv:2502.20257, eq:gauge_projs.
\end{definition}

\begin{lemma}[The placed Gauss projector is a projection]
    \label{thm:mpug_placed_gauss_projector_projection}
    \lean{MPOTensor.isStarProjection_gaussWindowProjector,
        MPOTensor.isStarProjection_placedGaussProjector}
    \leanok
    \discussion{7605}
    \uses{def:mpug_placed_gauss_projector}
    For every $N\geq2$ and every position $j$, the operator
    in~(\ref{eq:mpug_placed_gauss_projector}) satisfies
    \begin{align}
        P_j(R)^2       & =P_j(R), \\
        P_j(R)^\dagger & =P_j(R).
    \end{align}
\end{lemma}

\begin{proof}\leanok
    \uses{thm:mpug_gauss_projector_projection,
        thm:mpdo_cyclic_local_embedding_reindex}
    Let $U_\kappa$ be the permutation unitary determined by the coordinate
    bijection in Definition~\ref{def:mpug_placed_gauss_projector}, with the
    orientation
    \begin{align}
        U_\kappa\ket{((i_0,i_1),(a,b))}
         & =\ket{\kappa((i_0,i_1),(a,b))}.
        \notag
    \end{align}
    Simultaneous row and column reindexing is therefore the conjugation
    \begin{align}
        \widetilde P_G(R)
         & =U_\kappa P_G(R)U_\kappa^{-1}.
        \label{eq:mpug_gauss_reindex_conjugation}
    \end{align}
    Since $U_\kappa^{-1}=U_\kappa^\dagger$, Lemma~\ref{thm:mpug_gauss_projector_projection}
    and~(\ref{eq:mpug_gauss_reindex_conjugation}) give
    \begin{align}
        \widetilde P_G(R)^2
         & =U_\kappa P_G(R)^2U_\kappa^{-1}
        =\widetilde P_G(R),
        \label{eq:mpug_gauss_reindexed_idempotent} \\
        \widetilde P_G(R)^\dagger
         & =U_\kappa P_G(R)^\dagger U_\kappa^{-1}
        =\widetilde P_G(R).
        \label{eq:mpug_gauss_reindexed_adjoint}
    \end{align}
    Let $V_j$ be the permutation unitary that sends a chain configuration to
    its restriction to the cyclic window $(j,j+1)$ followed by its restriction
    to the complement. By Lemma~\ref{thm:mpdo_cyclic_local_embedding_reindex},
    \begin{align}
        V_jP_j(R)V_j^{-1}
         & =\widetilde P_G(R)\otimes\Id_{\mathrm{comp}}.
        \label{eq:mpug_gauss_window_tensor_identity}
    \end{align}
    Using $V_j^{-1}=V_j^\dagger$ and
    (\ref{eq:mpug_gauss_reindexed_idempotent})--(\ref{eq:mpug_gauss_reindexed_adjoint})
    in~(\ref{eq:mpug_gauss_window_tensor_identity}), we obtain
    \begin{align}
        V_jP_j(R)^2V_j^{-1}
         & =\bigl(\widetilde P_G(R)\otimes\Id_{\mathrm{comp}}\bigr)^2
        =\widetilde P_G(R)^2\otimes\Id_{\mathrm{comp}}
        =V_jP_j(R)V_j^{-1},                                                 \\
        V_jP_j(R)^\dagger V_j^{-1}
         & =\bigl(\widetilde P_G(R)\otimes\Id_{\mathrm{comp}}\bigr)^\dagger
        =\widetilde P_G(R)^\dagger\otimes\Id_{\mathrm{comp}}
        =V_jP_j(R)V_j^{-1}.
    \end{align}
    Conjugating by $V_j^{-1}$ gives $P_j(R)^2=P_j(R)$ and
    $P_j(R)^\dagger=P_j(R)$.
\end{proof}

\begin{definition}[Common fixed subspace]
    \label{def:mpug_common_fixed_subspace}
    \lean{LinearMap.commonFixedSubmodule}
    \leanok
    Let $\mathcal H$ be a complex vector space and let $(P_j)_{j\in J}$ be any
    family of endomorphisms of $\mathcal H$. Define
    their common fixed subspace by
    \begin{align}
        \mathcal V(P)
         & =\bigcap_{j\in J}\ker(P_j-\Id_{\mathcal H}).
    \end{align}
    No commutativity of the operators is required. For the averaged Gauss
    projectors, this is the common $+1$ eigenspace called the gauge-invariant
    subspace in
    \cite[lines~4327--4335 and line~5198]{FrancoRubio2025MPUGauging}.
    Here $P_j$ is an averaged projector, not an individual Gauss operator
    $\mathcal G_g$.
    % Source anchors: arXiv:2502.20257, lines 4327--4335 and line 5198.
\end{definition}

\begin{lemma}[Membership in the common fixed subspace]
    \label{thm:mpug_mem_common_fixed_subspace}
    \lean{LinearMap.mem_commonFixedSubmodule_iff}
    \leanok
    \uses{def:mpug_common_fixed_subspace}
    A vector $\psi\in\mathcal H$ belongs to $\mathcal V(P)$ if and only if
    $P_j\psi=\psi$ for every $j\in J$.
\end{lemma}

\begin{proof}\leanok
    \uses{def:mpug_common_fixed_subspace}
    Membership in an indexed intersection is equivalent to membership in
    every factor, and membership in $\ker(P_j-\Id_{\mathcal H})$ is equivalent
    to $P_j\psi=\psi$.
\end{proof}

\begin{lemma}[Nonzero-vector lower bound]
    \label{thm:mpug_common_fixed_subspace_finrank}
    \lean{LinearMap.one_le_finrank_commonFixedSubmodule}
    \leanok
    \uses{def:mpug_common_fixed_subspace}
    If $\mathcal H$ is finite-dimensional, $\psi\in\mathcal V(P)$, and
    $\psi\neq0$, then $1\leq\dim_{\C}\mathcal V(P)$.
\end{lemma}

\begin{proof}\leanok
    The common fixed subspace is nonzero because it contains $\psi$, so its
    dimension is at least one.
\end{proof}

\subsection{Defect completions and gauge-invariant subspaces}

\begin{definition}[Prescribed defect maps]
    \label{def:mpug_defect_maps}
    \lean{TNLean.Algebra.DefectMaps}
    \leanok
    \discussion{7569}
    Let $\mathcal H_{\mathrm m}$ be the matter space of two sites. Prescribed
    defect maps assign to every ordered pair $(a,b)\in G^2$ a defect subspace
    \begin{align}
        \mathcal K_{a,b} & \subseteq\mathcal H_{\mathrm m}
    \end{align}
    together with a map $D_{a,b}$ of $\mathcal H_{\mathrm m}$ whose
    restriction to that subspace is an isometry: for every
    $\ket{\xi}\in\mathcal K_{a,b}$,
    \begin{align}
        D_{a,b}^\dagger D_{a,b}\ket{\xi}
         & =\ket{\xi}.
        \label{eq:mpug_defect_isometry}
    \end{align}
    When the fusion scalars are not block independent, the state-level
    construction of~\cite[lines~4215--4326]{FrancoRubio2025MPUGauging}
    prescribes the modified fusion operators only on such subspaces.
    % Source anchor: arXiv:2502.20257, lemma:modif and the sentence closing
    % its proof.
\end{definition}

\begin{definition}[Full unitary completion class]
    \label{def:mpug_completion_class}
    \lean{TNLean.Algebra.DefectMaps.IsCompletion,
        TNLean.Algebra.DefectMaps.completionClass,
        TNLean.Algebra.DefectMaps.mem_completionClass_iff}
    \leanok
    \discussion{7569}
    \uses{def:mpug_defect_maps}
    The full completion class of prescribed defect maps
    $D=(D_{a,b})_{a,b\in G}$ is
    \begin{align}
        \mathfrak C(D)
         & =\left\{R=(U_{a,b})_{a,b\in G}:
        U_{a,b}\in\mathrm U(\mathcal H_{\mathrm m}),\
        U_{a,b}|_{\mathcal K_{a,b}}=D_{a,b}|_{\mathcal K_{a,b}}\right\}.
        \label{eq:mpug_completion_class}
    \end{align}
    As noted in~\cite[lines~4215--4326]{FrancoRubio2025MPUGauging}, when
    $\mathcal K_{a,b}$ is a proper subspace of the matter space, the unitary
    extension of $D_{a,b}$ is non-unique.
\end{definition}

\begin{lemma}[Existence of full unitary completions]
    \label{thm:mpug_completion_class_nonempty}
    \lean{TNLean.Algebra.DefectMaps.completionClass_nonempty,
        TNLean.Algebra.DefectMaps.exists_completion}
    \leanok
    \discussion{7702}
    \uses{def:mpug_defect_maps, def:mpug_completion_class}
    Let $\mathcal H_{\mathrm m}$ be finite-dimensional. Every family of
    prescribed defect maps satisfying (\ref{eq:mpug_defect_isometry}) admits a
    simultaneous unitary extension:
    \begin{align}
        \mathfrak C(D) & \neq\varnothing.
        \label{eq:mpug_completion_class_nonempty}
    \end{align}
    Equivalently, there is a family
    $R=(U_{a,b})_{(a,b)\in G^2}$ of unitary operators such that, for every
    $a,b\in G$ and every $\ket{\xi}\in\mathcal K_{a,b}$,
    \begin{align}
        U_{a,b}\ket{\xi}
         & =D_{a,b}\ket{\xi}.
        \label{eq:mpug_completion_exists}
    \end{align}
    This is the finite-dimensional extension step implicit in the freedom of
    the modified fusion operators at
    \cite[lines~4215--4326]{FrancoRubio2025MPUGauging}, where the subspaces
    $\mathcal K_{a,b}$ and their prescribed isometries enter as hypotheses.
    % Source anchor: arXiv:2502.20257, lemma:modif and its proof,
    % lines 4215--4326. The theorem does not assert that the defect data arise
    % from locally orthogonal MPS blocks or action tensors; their construction
    % remains in issues #7349, #7350, and #7569.
\end{lemma}

\begin{proof}\leanok
    \uses{def:mpug_defect_maps, def:mpug_completion_class}
    For each pair $(a,b)$, the restriction of $D_{a,b}$ to
    $\mathcal K_{a,b}$ is a linear isometry and therefore extends to a unitary
    operator on the whole finite-dimensional matter space. Choosing one such
    extension for every pair gives a family $R$ satisfying
    (\ref{eq:mpug_completion_exists}), hence a member of $\mathfrak C(D)$.
\end{proof}

\begin{lemma}[Comparison identity for a completion]
    \label{thm:mpug_completion_transport}
    \lean{TNLean.Algebra.DefectMaps.IsCompletion.mulVec_eq_prescribed,
        TNLean.Algebra.DefectMaps.IsCompletion.mulVec_eq_mulVec,
        TNLean.Algebra.DefectMaps.IsCompletion.unitaryFactorizationComparison_mulVec,
        TNLean.Algebra.DefectMaps.IsCompletion.unitaryFactorizationComparison_gaussLegAction_mulVec}
    \leanok
    \discussion{7569}
    \uses{def:mpug_completion_class,
        def:mpug_unitary_factorization_comparison,
        def:mpug_gauss_leg_action}
    Let $R=(U_{a,b})\in\mathfrak C(D)$, let
    $\ket{\xi}\in\mathcal K_{a,b}$ and $\ket{\eta}\in\mathcal K_{c,d}$, and
    assume the exact covariance
    \begin{align}
        D_{a,b}\ket{\xi} & =D_{c,d}\ket{\eta}.
        \label{eq:mpug_prescribed_covariance}
    \end{align}
    Then, as required for the modified Gauss laws
    of~\cite[lines~4325--4335]{FrancoRubio2025MPUGauging} when $(c,d)=T_g(a,b)$,
    \begin{align}
        U_{a,b}\ket{\xi}             & =D_{a,b}\ket{\xi},
        \label{eq:mpug_completion_agrees}                  \\
        U_{a,b}\ket{\xi}             & =U_{c,d}\ket{\eta},
        \label{eq:mpug_completion_covariance}              \\
        \lambda_{a,b}^{c,d}\ket{\xi} & =\ket{\eta}.
        \label{eq:mpug_completion_transport}
    \end{align}
\end{lemma}

\begin{proof}\leanok
    \uses{def:mpug_completion_class,
        thm:mpug_unitary_factorization_comparison_laws}
    Identity~(\ref{eq:mpug_completion_agrees}) is the defining property of a
    completion on the defect subspace. Applying it at $(a,b)$ and at $(c,d)$
    and inserting~(\ref{eq:mpug_prescribed_covariance})
    gives~(\ref{eq:mpug_completion_covariance}), an equality of vectors in the
    whole matter space. Multiplying that equality by $U_{c,d}^\dagger$ and
    using $U_{c,d}^\dagger U_{c,d}=\Id$ gives
    \begin{align}
        \lambda_{a,b}^{c,d}\ket{\xi}
         & =U_{c,d}^\dagger U_{a,b}\ket{\xi}
        =U_{c,d}^\dagger U_{c,d}\ket{\eta}
        =\ket{\eta}.
    \end{align}
    The adjoint is applied only after the equality in the whole matter space
    is available; no step asserts that $U_{c,d}^\dagger$ restricts to
    $D_{c,d}^\dagger$ on a defect subspace.
\end{proof}

\begin{definition}[Label action at one window]
    \label{def:mpug_chain_label_action}
    \lean{MPOTensor.chainLabelAction,
        MPOTensor.extractWindow_chainLabelAction_zero,
        MPOTensor.extractWindow_chainLabelAction_one}
    \leanok
    \discussion{7569}
    \uses{def:mpug_gauss_leg_action}
    Let $\alpha=(a_1,\ldots,a_N)$ be a configuration of gauge labels on a
    periodic chain of length $N\geq2$. For a window $(j,j+1)$ and $g\in G$,
    the configuration $T_{j,g}\alpha$ carries the labels
    \begin{align}
        \left(a_jg^{-1},ga_{j+1}\right)
    \end{align}
    on that window and agrees with $\alpha$ at every other site.
\end{definition}

\begin{lemma}[Composition of the label action]
    \label{thm:mpug_chain_label_action_mul}
    \lean{MPOTensor.chainLabelAction_one,
        MPOTensor.chainLabelAction_chainLabelAction,
        MPOTensor.chainLabelActionEquiv,
        MPOTensor.chainLabelActionEquiv_apply}
    \leanok
    \discussion{7569}
    \uses{def:mpug_chain_label_action}
    For all $g,h\in G$,
    \begin{align}
        T_{j,e}        & =\id,      \\
        T_{j,g}T_{j,h} & =T_{j,gh}.
        \label{eq:mpug_chain_label_action_mul}
    \end{align}
    In particular $T_{j,g}$ is a bijection of gauge-label configurations with
    inverse $T_{j,g^{-1}}$.
\end{lemma}

\begin{proof}\leanok
    \uses{def:mpug_chain_label_action}
    Outside the window both sides leave $\alpha$ unchanged. On the window,
    replacing the labels twice keeps only the second replacement, and the two
    labels are
    \begin{align}
        \left(a_jh^{-1}g^{-1},gha_{j+1}\right)
         & =\left(a_j(gh)^{-1},(gh)a_{j+1}\right).
    \end{align}
    At $g=e$ the labels are unchanged, so $T_{j,e}=\id$, and the composition
    law with $h=g^{-1}$ gives the inverse.
\end{proof}

\begin{definition}[Placed Gauss operator of a completion]
    \label{def:mpug_placed_gauss_operator}
    \lean{MPOTensor.gaussWindowOperator,
        MPOTensor.gaussWindowOperator_apply,
        MPOTensor.placedGaussOperator}
    \leanok
    \discussion{7569}
    \uses{def:mpug_gauss_operator,
        def:mpug_placed_gauss_projector,
        def:mpug_completion_class}
    Let $R\in\mathfrak C(D)$ and let a chain site carry one matter index and
    one gauge label. Relabelling the row and column coordinates of
    $\mathcal G_g$ by the bijection $\kappa$ of
    Definition~\ref{def:mpug_placed_gauss_projector} and embedding the result
    on the periodic window $(j,j+1)$ defines
    \begin{align}
        \widetilde{\mathcal G}^{j,j+1}_g(R)
         & =\operatorname{Embed}_{j,j+1}\bigl(\widetilde{\mathcal G}_g(R)\bigr),
        \label{eq:mpug_placed_gauss_operator}
    \end{align}
    the modified Gauss law
    of~\cite[lines~4325--4335]{FrancoRubio2025MPUGauging}, whose block from
    $(a,b)$ to $T_g(a,b)$ is $U_{T_g(a,b)}^\dagger U_{a,b}$.
\end{definition}

\begin{lemma}[The placed projector is the average of the placed Gauss operators]
    \label{thm:mpug_placed_gauss_projector_average}
    \lean{MPOTensor.placedGaussProjector_eq_average}
    \leanok
    \discussion{7569}
    \uses{def:mpug_placed_gauss_operator}
    For every window $(j,j+1)$,
    \begin{align}
        \widetilde P_j(R)
         & =\frac{1}{|G|}\sum_{g\in G}\widetilde{\mathcal G}^{j,j+1}_g(R).
        \label{eq:mpug_placed_gauss_projector_average}
    \end{align}
\end{lemma}

\begin{proof}\leanok
    \uses{def:mpug_gauss_projector, def:mpug_placed_gauss_operator}
    Comparing matrix entries, the relabelling is entrywise and the embedding
    multiplies each entry by the indicator that the two chain configurations
    agree outside the window. Both operations commute with the normalized sum
    over $G$ defining the local projector kernel
    in~(\ref{eq:mpug_gauss_projector}).
\end{proof}

\begin{definition}[Gauged vector of matter defect states]
    \label{def:mpug_gauged_vector}
    \lean{MPOTensor.gaugedMatter,
        MPOTensor.gaugedLabel,
        MPOTensor.gaugedVector,
        MPOTensor.gaugedVector_apply}
    \leanok
    \discussion{7569}
    Every chain configuration splits into its matter part and its gauge-label
    part. Given matter defect states $\ket{\Psi_\alpha}$ indexed by the
    gauge-label configurations $\alpha$ of a chain of length $N$, the gauged
    vector of \cite[lines~4325--4335]{FrancoRubio2025MPUGauging} is
    \begin{align}
        \ket{\Psi_{\mathrm{gauged}}}
         & =\sum_\alpha\ket{\Psi_\alpha}\otimes\ket{\alpha},
        \label{eq:mpug_gauged_vector}
    \end{align}
    whose component at a chain configuration is the component of
    $\ket{\Psi_\alpha}$ at its matter part, $\alpha$ being its gauge-label
    part.
\end{definition}

\begin{definition}[Local defect support and exact defect covariance]
    \label{def:mpug_defect_state_hypotheses}
    \lean{MPOTensor.HasDefectSupport,
        MPOTensor.HasPrescribedDefectCovariance}
    \leanok
    \discussion{7569}
    \uses{def:mpug_defect_maps, def:mpug_chain_label_action,
        def:mpug_gauged_vector}
    The matter defect states have local defect support when, for every
    $\alpha$ and every window $(j,j+1)$ whose adjacent labels are
    $(a_j,a_{j+1})$, the matter support of $\ket{\Psi_\alpha}$ lies in
    \begin{align}
        \mathcal K_{a_j,a_{j+1}}\otimes\mathcal H_{\mathrm{rest}}.
    \end{align}
    They satisfy exact defect covariance when
    \begin{align}
        \left(D_{a_j,a_{j+1}}\otimes\id_{\mathrm{rest}}\right)
        \ket{\Psi_\alpha}
         & =\left(D_{T_g(a_j,a_{j+1})}\otimes\id_{\mathrm{rest}}\right)
        \ket{\Psi_{T_{j,g}\alpha}}
        \label{eq:mpug_exact_defect_covariance}
    \end{align}
    for every $\alpha$, every window $(j,j+1)$, and every $g\in G$.
\end{definition}

\begin{lemma}[Universal invariant vector]
    \label{thm:mpug_universal_invariant_vector}
    \lean{MPOTensor.embedLocalOperator_mulVec_replaceWindow,
        MPOTensor.placedGaussOperator_mulVec_gaugedVector,
        MPOTensor.placedGaussProjector_mulVec_gaugedVector}
    \leanok
    \discussion{7569}
    \uses{def:mpug_placed_gauss_operator,
        def:mpug_placed_gauss_projector,
        def:mpug_completion_class,
        def:mpug_gauged_vector,
        def:mpug_defect_state_hypotheses}
    Assume local defect support and exact defect
    covariance~(\ref{eq:mpug_exact_defect_covariance}), which as recorded
    in~\cite{gap:fbc25_state_level_gauging_covariance} enter here as hypotheses
    rather than consequences of locally orthogonal MPS blocks. For every
    $R\in\mathfrak C(D)$, every window $(j,j+1)$, and every $g\in G$,
    \begin{align}
        \widetilde{\mathcal G}^{j,j+1}_g(R)\ket{\Psi_{\mathrm{gauged}}}
         & =\ket{\Psi_{\mathrm{gauged}}},
        \label{eq:mpug_universal_invariant_vector} \\
        \widetilde P_j(R)\ket{\Psi_{\mathrm{gauged}}}
         & =\ket{\Psi_{\mathrm{gauged}}}.
        \label{eq:mpug_universal_invariant_proj}
    \end{align}
\end{lemma}

\begin{center}
    % Formula: \widetilde{\mathcal G}^{j,j+1}_g(R)\ket{\Psi_{\mathrm{gauged}}}
    % =\ket{\Psi_{\mathrm{gauged}}}, equation
    % (\ref{eq:mpug_universal_invariant_vector}); source: arXiv:2502.20257,
    % lines 4325--4335.
    % Ink-to-index map: the wide lower box is the gauged vector, one leg per
    % chain site carrying the matter index and the gauge label of that site;
    % sites 1, j, j+1, N are drawn.  The upper box is the placed Gauss
    % operator acting on the window (j,j+1) and as the identity elsewhere.
    % Contraction set: the two window legs between the operator and the
    % vector; the output legs remain open.
    % Boundary signature: (virtual-west=0 | virtual-east=0 |
    % physical-up=4 | physical-down=0) on both sides.
    \tenkzkernel
    \begin{tenkz}[rows={op,ket}, cols=4, bonds=none]
        \tn[at=(1,2), name=G, skin=box, wide=2, ports={90@1:physical:$j$, 90@2:physical:$j{+}1$, 270@1:physical, 270@2:physical}]{\widetilde{\mathcal G}^{j,j+1}_g(R)}
        \tn[at=(2,1), name=P, skin=box, wide=4, ports={90@1:physical:$1$, 90@2:physical, 90@3:physical, 90@4:physical:$N$}]{\Psi_{\mathrm{gauged}}}
        \tn[at=(1,1), void=sealed]{}
        \tn[at=(1,4), void=sealed]{}
        \tnwire{G.270@1}{P.90@2}
        \tnwire{G.270@2}{P.90@3}
    \end{tenkz}
    \quad $=$ \quad
    \begin{tenkz}[rows={ket}, cols=4, bonds=none]
        \tn[at=(1,1), name=P, skin=box, wide=4, ports={90@1:physical:$1$, 90@2:physical:$j$, 90@3:physical:$j{+}1$, 90@4:physical:$N$}]{\Psi_{\mathrm{gauged}}}
    \end{tenkz}
\end{center}

\begin{proof}\leanok
    \uses{thm:mpug_completion_transport,
        thm:mpug_chain_label_action_mul,
        thm:mpug_placed_gauss_projector_average,
        def:mpug_gauss_operator}
    Fix an output chain configuration and let $\alpha$ be its gauge-label part,
    with $(a,b)=(a_j,a_{j+1})$ the two labels adjacent to the window. Among the
    columns contributing to this fixed output row of
    $\widetilde{\mathcal G}^{j,j+1}_g(R)$, only the gauge labels
    $T_{j,g^{-1}}\alpha$ survive. Thus, with
    $\beta=T_{j,g^{-1}}\alpha$, the component of
    $\widetilde{\mathcal G}^{j,j+1}_g(R)\ket{\Psi_{\mathrm{gauged}}}$ equals
    the component of
    \begin{align}
        \left(U_{a,b}^\dagger U_{ag,g^{-1}b}
        \otimes\id_{\mathrm{rest}}\right)\ket{\Psi_\beta},
    \end{align}
    since the labels of $\beta$ adjacent to the window are $(ag,g^{-1}b)$ and
    $T_g(ag,g^{-1}b)=(a,b)$, whence $T_{j,g}\beta=\alpha$ by the composition
    law~(\ref{eq:mpug_chain_label_action_mul}). Local defect support puts the
    matter slices of $\ket{\Psi_\beta}$ and $\ket{\Psi_\alpha}$ over that
    window in $\mathcal K_{ag,g^{-1}b}$ and $\mathcal K_{a,b}$, and exact
    covariance~(\ref{eq:mpug_exact_defect_covariance}) relates them, so the
    comparison identity~(\ref{eq:mpug_completion_transport}) turns that vector
    into $\ket{\Psi_\alpha}$. Identifying the label configuration $\beta$ with
    $\alpha$ through the bijection $T_{j,g}$ is the reindexing of the sum
    over gauge-label configurations in~(\ref{eq:mpug_gauged_vector}). Averaging~(\ref{eq:mpug_universal_invariant_vector})
    over $g$ via~(\ref{eq:mpug_placed_gauss_projector_average}) yields~(\ref{eq:mpug_universal_invariant_proj}).
\end{proof}

\begin{definition}[Gauge-invariant subspace of a completion]
    \label{def:mpug_gauge_invariant_subspace}
    \lean{MPOTensor.gaugeInvariantSubspace}
    \leanok
    \discussion{7569}
    \uses{def:mpug_common_fixed_subspace,
        def:mpug_completion_class,
        def:mpug_placed_gauss_projector}
    For $R\in\mathfrak C(D)$ and a chain of length $N\geq2$, the
    gauge-invariant subspace is the common fixed subspace
    \begin{align}
        \mathcal V_N(R)
         & =\left\{\ket{\phi}:
        \widetilde P_j(R)\ket{\phi}=\ket{\phi}
        \text{ for every }j\right\}.
        \label{eq:mpug_gauge_invariant_subspace}
    \end{align}
    As noted in~\cite[lines~4325--4335]{FrancoRubio2025MPUGauging}, the
    modified projectors on neighbouring windows need not commute.
\end{definition}

\begin{lemma}[The gauged vector spans at least one dimension]
    \label{thm:mpug_gauge_invariant_subspace_finrank}
    \lean{MPOTensor.gaugedVector_mem_gaugeInvariantSubspace,
        MPOTensor.one_le_finrank_gaugeInvariantSubspace}
    \leanok
    \discussion{7569}
    \uses{def:mpug_gauge_invariant_subspace,
        def:mpug_completion_class,
        def:mpug_gauged_vector,
        def:mpug_defect_state_hypotheses}
    Under local defect support and exact defect
    covariance~(\ref{eq:mpug_exact_defect_covariance}),
    \begin{align}
        \ket{\Psi_{\mathrm{gauged}}} & \in\mathcal V_N(R)
    \end{align}
    for every $R\in\mathfrak C(D)$. If $\ket{\Psi_{\mathrm{gauged}}}\neq0$, then
    \begin{align}
        \dim_{\C}\mathcal V_N(R) & \geq1.
    \end{align}
    Whereas \cite[line~5198]{FrancoRubio2025MPUGauging} asks how
    $\dim_{\C}\mathcal V_N(R)$ grows with $N$, neither the minimum nor the
    maximum of $\dim_{\C}\mathcal V_N(R)$ over $\mathfrak C(D)$ is asserted
    here.
\end{lemma}

\begin{proof}\leanok
    \uses{thm:mpug_universal_invariant_vector,
        thm:mpug_mem_common_fixed_subspace,
        thm:mpug_common_fixed_subspace_finrank}
    Membership is the projector identity in
    Lemma~\ref{thm:mpug_universal_invariant_vector} at every window. A
    nonzero member of a subspace forces its dimension to be at least one.
\end{proof}
